% Lösungen Einheit 5 von 5 – Besondere Linien im Dreieck und Satz des Thales (zusammenbau v0.1)
\einheitenkopf{Einheit 5 von 5 · Besondere Linien im Dreieck und Satz des Thales}

%% TODO Lösung der Erklärzeile 23g) fehlt in der Bank
\erg{23}{a) Höhe \quad b) Kreisbögen um A und um B mit gleichem Radius (größer als die halbe Strecke); Gerade durch die beiden Schnittpunkte; sie trifft AB in der Mitte M(5|3) \quad c) Kreisbögen um A und um B mit gleichem Radius (größer als die halbe Strecke); Gerade durch die beiden Schnittpunkte; sie trifft AB in der Mitte M(4|3) \quad d) Kreisbögen um A und um B mit gleichem Radius (größer als die halbe Strecke); Gerade durch die beiden Schnittpunkte; sie trifft AB in der Mitte M(6|3) \quad e) Kreisbögen um A und um B mit gleichem Radius (größer als die halbe Strecke); Gerade durch die beiden Schnittpunkte; sie trifft AB in der Mitte M(4|2) \quad f) Kreisbögen um A und um B mit gleichem Radius (größer als die halbe Strecke); Gerade durch die beiden Schnittpunkte; sie trifft AB in der Mitte M(4|3) \quad g) \ldots \quad h) Die Mittelsenkrechten schneiden sich in U(4|2); Kreis um U durch A, B und C, Radius ≈ 3,2 Einheiten \quad i) Kreisbogen um den Scheitel schneidet beide Schenkel; um diese Punkte gleich große Bögen; Strahl vom Scheitel durch deren Schnittpunkt. Jeder Teilwinkel: $64^\circ : 2 = 32^\circ$ \quad j) Die Winkelhalbierenden schneiden sich in I(3|3); das Lot von I auf eine Seite ist der Radius ≈ 2 Einheiten; Kreis um I \quad k) Lot von C auf die Gerade AB, Fußpunkt D(4|1) \quad l) Jede Seitenmitte mit der gegenüberliegenden Ecke verbinden; Schnittpunkt S(4|3) \quad m) Halbkreis um die Mitte von AB mit dem Radius $8 : 2 = 4$ cm; der rechte Winkel liegt bei C \quad n) AB ist ein Durchmesser, und C liegt auf dem Kreis, also auf dem Thaleskreis über AB; nach dem Satz des Thales ist der Winkel bei C ein rechter. \quad o) D liegt auf dem Kreis um M durch A und C; AC ist ein Durchmesser dieses Kreises. Nach dem Satz des Thales ist der Winkel bei D ein rechter.}
\erg{24}{a) Fehler: Das ist die Seitenhalbierende. Die Höhe ist das Lot von C auf AB und steht senkrecht auf AB. \quad b) Jeder Punkt der Mittelsenkrechten von AB ist von A und B gleich weit entfernt, jeder Punkt der Mittelsenkrechten von BC von B und C. Der Schnittpunkt liegt auf beiden, also ist er von A, B und C gleich weit entfernt. \quad c) Schnittpunkt zweier Mittelsenkrechten U(5|4); Entfernung = 5 km}
